\chapter*{Resumen}
\addcontentsline{toc}{chapter}{Resumen}

Este Trabajo Fin de Máster estudia IA agéntica para personalizar la gestión
conversacional de citas sanitarias en un entorno simulado. La arquitectura
principal, M5, distribuye cinco responsabilidades LLM: contexto, perfiles,
cumplimiento, personalización y orquestación. Se compara con S1, un agente
único, y H, un flujo híbrido programado. Las tres variantes comparten
herramientas y controles de acceso, confirmación y escritura.

El prototipo utiliza perfiles y agendas sintéticos y un subconjunto FHIR R4
para reservar, consultar, reprogramar y cancelar citas y configurar
recordatorios sin envío externo. La personalización adapta respuestas y
ordena huecos compatibles, sin modificar permisos ni resultados. CitaLab
permite interactuar, inspeccionar llamadas y revisar conversaciones.

La evaluación ejecuta 276 episodios con dos modelos y un usuario simulado,
separados del corpus empleado para ajustar las instrucciones. El éxito se
comprueba sobre la agenda, y la calidad de la conversación se valora con un
jurado de dos LLM, contrastado con una revisión humana ciega de una muestra y
con el oráculo. No se observa una ventaja significativa de M5 en éxito, mientras que
su coste es sistemáticamente mayor: el doble de llamadas y unas 2,4 veces los
tokens de H, que obtiene una calidad igual o superior. No se observan
escrituras sin confirmación ni cambios en citas ajenas, pero confirmar
correctamente una operación no garantiza cumplir la intención de la persona. Las conclusiones se limitan al
entorno sintético estudiado. El código, los
datos y los resultados están disponibles en \url{https://github.com/Andros5/citalab}.

\vspace{1em}
\noindent\textbf{Palabras clave:} IA agéntica, personalización, sistemas
multiagente, citas sanitarias, FHIR, evaluación.
